% =====================================================================
% tesi/bibliografia.tex - GENERATO, non modificare a mano
% =====================================================================
% Prodotto da tools/build-bibliography.py dalla tabella FONTI di
% tools/build-source-map.py, che è la stessa da cui nascono le note di
% docs/fonti/. Modificare quella tabella e rigenerare: una correzione fatta
% qui sparisce alla corsa successiva.
%
% La chiave di ogni voce è lo slug della fonte, quindi \cite{pokered} qui e
% docs/fonti/pokered.md nel vault sono la stessa fonte vista da due parti.

\begin{thebibliography}{999}
\addcontentsline{toc}{chapter}{Bibliografia}
\raggedright
\setlength{\labelsep}{0.5em}
\setlength{\itemsep}{0.8ex}

\bibitem{pokered}
  \textbf{pret/pokered}.
  \\ \url{https://github.com/pret/pokered}
  \\ \emph{Livello 1, disassemblato o documentazione dell'hardware; letta. Track: BRI.}
  \\ Disassemblaggio completo di Pokemon Rosso e Blu che ricompila in una ROM identica all'originale. Contiene le macro delle strutture dati, la mappa della memoria di lavoro, le costanti del protocollo seriale e il codice del Centro Scambi.

\bibitem{pokecrystal}
  \textbf{pret/pokecrystal}.
  \\ \url{https://github.com/pret/pokecrystal}
  \\ \emph{Livello 1, disassemblato o documentazione dell'hardware; letta. Track: BRI.}
  \\ Disassemblaggio di Pokemon Cristallo. Oltre alle strutture di generazione 2 contiene la routine di calcolo delle statistiche, la tabella dei caratteri e le due strutture di invio sul cavo, quella nativa e quella del Time Capsule.

\bibitem{pokeemerald}
  \textbf{pret/pokeemerald}.
  \\ \url{https://github.com/pret/pokeemerald}
  \\ \emph{Livello 1, disassemblato o documentazione dell'hardware; letta. Track: BRI, SME.}
  \\ Decompilazione di Pokemon Smeraldo in C e assembly ARM che ricompila in una ROM identica. Contiene la struttura cifrata del Pokemon, il calcolo del checksum, la mappa dei settori del salvataggio, la chiave di cifratura e la maschera sulle quantità degli oggetti.

\bibitem{pokefirered}
  \textbf{pret/pokefirered}.
  \\ \url{https://github.com/pret/pokefirered}
  \\ \emph{Livello 1, disassemblato o documentazione dell'hardware; letta. Track: BRI, LDN, SME.}
  \\ Decompilazione di Rosso Fuoco e Verde Foglia. Struttura del salvataggio diversa da quella di Smeraldo in ogni offset che conta: chiave di cifratura, conteggio della squadra, denaro e tasche dello zaino.

\bibitem{pokeruby}
  \textbf{pret/pokeruby}.
  \\ \url{https://github.com/pret/pokeruby}
  \\ \emph{Livello 1, disassemblato o documentazione dell'hardware; letta. Track: BRI, SME.}
  \\ Decompilazione di Rubino e Zaffiro. Non ha alcuna chiave di cifratura, quindi le quantità degli oggetti sono in chiaro, ed è l'unico dei tre giochi Gen 3 a comportarsi così.

\bibitem{pandocs}
  \textbf{Pan Docs, trasferimento seriale}.
  \\ \url{https://gbdev.io/pandocs/Serial_Data_Transfer_(Link_Cable).html}
  \\ \emph{Livello 1, disassemblato o documentazione dell'hardware; letta. Track: BRI.}
  \\ Riferimento tecnico dell'hardware seriale del Game Boy: i registri SB e SC con i loro bit, le frequenze del clock interno, e il comportamento del clock esterno, che il gioco accetta a qualunque velocità fino a 500 kHz e anche a intervalli irregolari.

\bibitem{gbatek}
  \textbf{GBATEK, multiboot}.
  \\ \url{https://problemkaputt.de/gbatek-bios-multi-boot-single-game-pak.htm}
  \\ \emph{Livello 1, disassemblato o documentazione dell'hardware; letta. Track: BRI.}
  \\ Specifica del multiboot del Game Boy Advance: la stretta di mano con 0x6200 e la risposta 0x0000, le dimensioni ammesse del programma, gli indirizzi di destinazione in memoria di lavoro, la cifratura del trasferimento e i vincoli di temporizzazione.

\bibitem{copetti}
  \textbf{Copetti, architettura del GBA}.
  \\ \url{https://www.copetti.org/writings/consoles/game-boy-advance/}
  \\ \emph{Livello 1, disassemblato o documentazione dell'hardware; letta. Track: BRI.}
  \\ Descrizione architetturale del Game Boy Advance: le due memorie di lavoro con le loro dimensioni e velocità, il bus della cartuccia, il buffer di prefetch e il meccanismo della retrocompatibilità con il Game Boy.

\bibitem{devlog-ptgb}
  \textbf{Dev log di Poke Transporter GB}.
  \\ \url{https://www.austinthomasweber.com/poke-transporter-gb}
  \\ \emph{Livello 4, articolo, blog o ricerca applicata; letta. Track: BRI.}
  \\ Serie di undici articoli in cui l'autore del ponte di riferimento racconta il proprio processo, dall'architettura della console al formato di salvataggio, dal protocollo del cavo alla scoperta dell'exploit, fino all'iniezione dell'evento e alla gestione della grafica.

\bibitem{ptgb}
  \textbf{Poke Transporter GB}.
  \\ \url{https://github.com/Striaton-Lab-Team/Poke_Transporter_GB}
  \\ \emph{Livello 3, implementazione di riferimento; letta. Track: BRI.}
  \\ Il ponte di riferimento: homebrew per Game Boy Advance in multiboot che trasferisce Pokemon da Gen 1 e 2 a Gen 3 usando il cavo del Game Boy Color, sotto licenza MIT. Contiene un generatore di payload con assemblatore Z80 proprio e tabelle di valori di ROM per lingua.

\bibitem{pccs}
  \textbf{PCCS}.
  \\ \url{https://github.com/Striaton-Lab-Team/Pokemon-Community-Conversion-Standard}
  \\ \emph{Livello 3, implementazione di riferimento; letta. Track: BRI.}
  \\ Libreria C++ che definisce come convertire un Pokemon da Gen 1 e 2 a Gen 3. Il README documenta quattro metodi campo per campo; il codice della release corrente implementa il comportamento di uno solo.

\bibitem{gen3togenx}
  \textbf{Pokemon-Gen3-to-Gen-X}.
  \\ \url{https://github.com/Lorenzooone/Pokemon-Gen3-to-Gen-X}
  \\ \emph{Livello 3, implementazione di riferimento; letta. Track: BRI.}
  \\ Homebrew per Game Boy Advance che scambia Pokemon fra Gen 3 e Gen 1 e 2 usando il protocollo normale e non exploit, con in più la gestione dell'orologio interno di Rubino, Zaffiro e Smeraldo.

\bibitem{pokemongb-online}
  \textbf{PokemonGB\_Online\_Trades}.
  \\ \url{https://github.com/Lorenzooone/PokemonGB_Online_Trades}
  \\ \emph{Livello 3, implementazione di riferimento; letta. Track: BRI.}
  \\ Implementazione in Python del protocollo di scambio di Gen 1, 2 e 3, funzionante sia su un adattatore USB per il cavo sia sull'emulatore BGB, con multiboot per il lato Gen 3.

\bibitem{cable-link}
  \textbf{CableClub/cable-link}.
  \\ \url{https://github.com/CableClub/cable-link}
  \\ \emph{Livello 3, implementazione di riferimento; letta. Track: BRI.}
  \\ Circuito stampato in formato KiCad, con gerber pronti alla produzione, e firmware per Raspberry Pi Pico che parla il protocollo di scambio di generazione 1. Apache 2.0, del 2021.

\bibitem{pksploit}
  \textbf{PkSploit}.
  \\ \url{https://github.com/binarycounter/PkSploit}
  \\ \emph{Livello 3, implementazione di riferimento; letta. Track: BRI, SME.}
  \\ Suite che dumpa ROM e salvataggio e riscrive il salvataggio di qualunque cartuccia Game Boy e Game Boy Color, usando come vettore una cartuccia di generazione 1, un cavo Link e un microcontrollore compatibile Arduino.

\bibitem{cableclubhack}
  \textbf{vaguilar/pokemon-red-cable-club-hack}.
  \\ \url{https://github.com/vaguilar/pokemon-red-cable-club-hack}
  \\ \emph{Livello 3, implementazione di riferimento; letta. Track: BRI.}
  \\ Implementazione originale dell'exploit del Centro Scambi con Arduino e Python. Documenta il cablaggio dei pin, cioè uscita seriale sul pin 6, ingresso sul 3, clock sul 2 e massa, e contiene il payload in un array dentro un'intestazione Arduino.

\bibitem{linkhack}
  \textbf{Phasip/PokemonLinkHack}.
  \\ \url{https://github.com/Phasip/PokemonLinkHack}
  \\ \emph{Livello 3, implementazione di riferimento; letta. Track: BRI.}
  \\ Catena di exploit che dal buffer overflow del cavo arriva a installare programmi persistenti dentro il deposito Pokemon: un selettore nel primo slot e i programmi nei successivi, sopravvivendo ai riavvii grazie al salvataggio della cartuccia.

\bibitem{blog-phasip}
  \textbf{Blog di Phasip}.
  \\ \url{https://www.sn1.se/posts/pokemon/}
  \\ \emph{Livello 4, articolo, blog o ricerca applicata; letta. Track: BRI.}
  \\ Racconto della catena di exploit con i numeri: 198 byte di payload utile, i programmi che risiedono negli slot del deposito, e il limite dichiarato che fuori dal buffer si possono scrivere solo nomi di Pokemon predefiniti.

\bibitem{arduino-poke-gen2}
  \textbf{arduino-poke-gen2}.
  \\ \url{https://github.com/stevenchaulk/arduino-poke-gen2}
  \\ \emph{Livello 3, implementazione di riferimento; letta. Track: BRI.}
  \\ Adattamento a generazione 2 di una macchina a stati Arduino per lo scambio, provato su Cristallo, con schema di cablaggio incluso.

\bibitem{arduino-boy}
  \textbf{pepijndevos/arduino-boy}.
  \\ \url{https://github.com/pepijndevos/arduino-boy}
  \\ \emph{Livello 3, implementazione di riferimento; letta. Track: BRI.}
  \\ Deposito Pokemon su Arduino per i giochi di generazione 1: memorizza un Pokemon nella memoria non volatile della scheda e lo scambia con la console. Il codice del protocollo deriva a sua volta da un progetto precedente chiamato gameboy-spoof.

\bibitem{gba-link-connection}
  \textbf{gba-link-connection}.
  \\ \url{https://github.com/afska/gba-link-connection}
  \\ \emph{Livello 3, implementazione di riferimento; letta. Track: BRI.}
  \\ Libreria C++ per il porto seriale del Game Boy Advance, con moduli separati per la modalità multiplayer, l'invio di multiboot ad altre console, l'adattatore wireless, il protocollo Joybus verso Wii e GameCube, le carte e-Reader e il Mobile Adapter GB.

\bibitem{reon}
  \textbf{Progetto REON}.
  \\ \url{https://github.com/REONTeam}
  \\ \emph{Livello 3, implementazione di riferimento; letta. Track: BRI.}
  \\ Ricostruzione dell'infrastruttura di rete del Mobile Adapter GB, l'accessorio che dava funzioni online ai giochi Game Boy. Comprende una libreria del protocollo in C, un server, un emulatore per BGB, un adattatore su Arduino e una utilità dedicata allo scambio nel Trade Corner.

\bibitem{gen3distributions}
  \textbf{Goppier/GEN3PokemonDistributions}.
  \\ \url{https://github.com/Goppier/GEN3PokemonDistributions}
  \\ \emph{Livello 3, implementazione di riferimento; letta. Track: BRI.}
  \\ Raccolta di cartucce di distribuzione per eventi di generazione 3, con la procedura operativa per usarle fra due console.

\bibitem{stadium-ace}
  \textbf{MrCheeze/pokestadium-ace}.
  \\ \url{https://github.com/MrCheeze/pokestadium-ace}
  \\ \emph{Livello 3, implementazione di riferimento; letta. Track: BRI.}
  \\ Esecuzione di codice arbitrario su Pokemon Stadium per Nintendo 64, che passa dal sistema di scambio e presuppone di avere già il controllo su un gioco di generazione 1. Documenta gli indirizzi dove i box vengono convertiti dal formato Gen 1 al formato Stadium.

\bibitem{usb-gba-multiboot}
  \textbf{usb-gba-multiboot}.
  \\ \url{https://github.com/tangrs/usb-gba-multiboot}
  \\ \emph{Livello 3, implementazione di riferimento; letta. Track: BRI.}
  \\ Firmware per microcontrollore Teensy e software su PC per caricare un programma nella console via USB, sfruttando il multiboot.

\bibitem{rom-sender}
  \textbf{gba-link-cable-rom-sender}.
  \\ \url{https://github.com/FIX94/gba-link-cable-rom-sender}
  \\ \emph{Livello 3, implementazione di riferimento; letta. Track: BRI.}
  \\ Homebrew per GameCube e Wii che invia un programma multiboot alla console collegata, leggendo i file da una cartella sulla scheda di memoria.

\bibitem{kinnay-ldn}
  \textbf{kinnay, protocollo LDN}.
  \\ \url{https://github.com/kinnay/NintendoClients/wiki/LDN-Protocol}
  \\ \emph{Livello 1, disassemblato o documentazione dell'hardware; letta. Track: LDN.}
  \\ Specifica del protocollo di rete locale della Nintendo Switch: action frame proprietario trasmesso ogni cento millisecondi, canali radio usati, struttura dell'annuncio campo per campo, tre livelli di cifratura con la derivazione delle chiavi, e sequenza di connessione con assegnazione degli indirizzi.

\bibitem{frlg-ldn-trade}
  \textbf{frlg-ldn-trade}.
  \\ \url{https://github.com/unlimitedcoder2/frlg-ldn-trade}
  \\ \emph{Livello 3, implementazione di riferimento; letta. Track: LDN.}
  \\ Proof of concept che fa scambiare Pokemon a un PC con Rosso Fuoco e Verde Foglia in esecuzione su Switch o Switch 2, simulando un giocatore che si collega come capo sessione. Richiede Linux, Python 3.12 o successivo, le chiavi della console, almeno due strutture .pk3 e un gioco portato avanti fino allo sbloccio della sala degli scambi. Licenza AGPLv3.

\bibitem{ldnd}
  \textbf{unlimitedcoder2/ldnd}.
  \\ \url{https://github.com/unlimitedcoder2/ldnd}
  \\ \emph{Livello 1, disassemblato o documentazione dell'hardware; letta. Track: LDN.}
  \\ Demone in C, licenza GPL-2.0, che porta lo stack wireless di Linux su Windows: collega il kernel Linux come libreria statica tramite LKL dentro un eseguibile costruito con MinGW, riceve l adattatore USB attraverso WinUSB e gli fa caricare i driver e i file di linux-firmware. Espone il servizio su una pipe con nome, e la sua riga di comando passa parametri di modulo del kernel fra cui rtw88\_usb.switch\_usb\_mode=0.

\bibitem{pmr-discord}
  \textbf{Pokemon Multiplayer Research, canali di supporto}.
  \\ \url{https://discord.gg/nBnTrv3UMn}
  \\ \emph{Livello 4, articolo, blog o ricerca applicata; letta. Track: LDN, BRI.}
  \\ Server della community che sviluppa lo scambio via rete locale sui giochi di generazione 3 su Switch. I canali di supporto e generale contengono la casistica reale degli adattatori wireless, i sintomi dei guasti con le loro cause, e le dichiarazioni degli autori sui limiti del progetto.

\bibitem{ldn-mitm}
  \textbf{ldn\_mitm}.
  \\ \url{https://github.com/spacemeowx2/ldn_mitm}
  \\ \emph{Livello 3, implementazione di riferimento; letta. Track: LDN.}
  \\ Modulo di sistema per Switch che sostituisce il servizio di rete locale ed emula la scansione delle console vicine usando la rete locale via UDP.

\bibitem{switch-lan-play}
  \textbf{switch-lan-play}.
  \\ \url{https://github.com/spacemeowx2/switch-lan-play}
  \\ \emph{Livello 3, implementazione di riferimento; letta. Track: LDN.}
  \\ Client e server che intercettano il traffico di rete della console con libpcap e lo incapsulano in un protocollo minimale, per far credere alle console di essere sulla stessa rete locale.

\bibitem{bulbapedia}
  \textbf{Bulbapedia, pagine sui formati}.
  \\ \url{https://bulbapedia.bulbagarden.net/wiki/Pok\%C3\%A9mon_data_structure_(Generation_III)}
  \\ \emph{Livello 2, wiki o riferimento di dominio; letta. Track: BRI, SME, LDN.}
  \\ Insieme di pagine enciclopediche sui formati dei dati e dei salvataggi delle tre generazioni, sui valori individuali, sul valore di personalità, sulla codifica dei caratteri e sugli indici di specie.

\bibitem{gcri-discord}
  \textbf{Glitch City Research Institute, canali}.
  \\ \url{https://discord.com/invite/EA7jxJ6}
  \\ \emph{Livello 4, articolo, blog o ricerca applicata; letta. Track: BRI.}
  \\ Server della community che studia i difetti sfruttabili dei giochi Pokemon. I canali per generazione contengono il lavoro corrente, che il wiki recepisce con ritardo o non recepisce affatto.

\bibitem{glitchcity}
  \textbf{Glitch City Wiki}.
  \\ \url{https://glitchcity.wiki}
  \\ \emph{Livello 2, wiki o riferimento di dominio; letta. Track: BRI.}
  \\ Catalogo della ricerca sui glitch dei giochi di generazione 1 e 2, compresi i metodi di esecuzione di codice raggiungibili giocando e quelli che passano dal cavo.

\bibitem{pokemon-automation}
  \textbf{Pokemon Automation}.
  \\ \url{https://pokemonautomation.github.io/}
  \\ \emph{Livello 3, implementazione di riferimento; letta. Track: AUT, LDN.}
  \\ Progetto che automatizza le parti ripetitive dei giochi Pokemon su Nintendo Switch con oltre cento programmi. È un anello di controllo chiuso su un sistema che non espone stato: percepisce il fotogramma video e in alcuni titoli l'audio, attua tramite un controller emulato da un microcontrollore, e decide con un programma per compito. Il perimetro dichiarato è console non modificate e nessun accesso alla memoria.

\bibitem{gambatte-gamelink}
  \textbf{Gambatte con GameLink su TCP}.
  \\ \url{https://gbatemp.net/threads/mission-wireless-trading-on-gen1-and-gen2-pokemon-games.632492/}
  \\ \emph{Livello 5, forum o community; letta. Track: BRI.}
  \\ Discussione in cui un membro compila una versione di Gambatte con il collegamento seriale del Game Boy emulato su TCP, funzionante sia come client sia come server, e riferisce di aver scambiato e combattuto nei giochi di generazione 1 e 2 fra dispositivi diversi, Switch compreso.

\bibitem{gbatemp-vc-save}
  \textbf{GBAtemp, correzione dei salvataggi in Virtual Console}.
  \\ \url{https://gbatemp.net/threads/tutorial-fix-all-save-problems-for-pokemon-games-vc-gba.433266/}
  \\ \emph{Livello 4, articolo, blog o ricerca applicata; letta. Track: 3DS.}
  \\ Tutorial del 2016 che corregge l impossibilità di salvare nei giochi Pokemon per Game Boy Advance iniettati come Virtual Console su Nintendo 3DS, tramite modifica esadecimale della ROM, e in una seconda parte rimuove il messaggio di salvataggio corrotto con offset specifici per versione e lingua.

\bibitem{projectpokemon}
  \textbf{Project Pokemon, discussioni}.
  \\ \url{https://projectpokemon.org/home/forums/topic/64794-pokemon-emerald-items-are-in-the-right-bag-using-the-app-but-when-i-load-it-into-a-cartridge-they-go-in-the-wrong-slots/}
  \\ \emph{Livello 5, forum o community; letta. Track: SME, BRI.}
  \\ Tre discussioni lette: una su un dispositivo che fa da sorgente di clock per il protocollo e ha completato scambi via internet, una su un salvataggio assente invece che corrotto su cartuccia contraffatta, e una su un editor che ha identificato male il gioco facendo finire gli oggetti negli slot sbagliati.

\bibitem{pokeyellow}
  \textbf{pret/pokeyellow}.
  \\ \url{https://github.com/pret/pokeyellow}
  \\ \emph{Livello 1, disassemblato o documentazione dell'hardware; letta. Track: BRI.}
  \\ Disassemblaggio di Pokemon Giallo, clonato e confrontato campo per campo con quello di Rosso e Blu.

\bibitem{pokegold}
  \textbf{pret/pokegold}.
  \\ \url{https://github.com/pret/pokegold}
  \\ \emph{Livello 1, disassemblato o documentazione dell'hardware; letta. Track: BRI.}
  \\ Disassemblaggio di Pokemon Oro e Argento, clonato e confrontato con quello di Cristallo.

\bibitem{video-goppier}
  \textbf{Goppier, i due video sul primo ponte}.
  \\ \url{https://www.youtube.com/watch?v=Qcp4vxyaUJc}
  \\ \emph{Livello 4, articolo, blog o ricerca applicata; letta. Track: BRI.}
  \\ Le due sole testimonianze esistenti sul primo ponte fra generazioni: l'aggiornamento di sviluppo, trascritto e letto per intero, e il video precedente che presenta il dispositivo, privo di sottotitoli di qualunque tipo e quindi ancora da trascrivere.

\bibitem{video-trascritti}
  \textbf{Quattro video trascritti su formati, cavo e strumenti}.
  \\ \url{https://www.youtube.com/watch?v=VVbRe7wr3G4}
  \\ \emph{Livello 4, articolo, blog o ricerca applicata; letta. Track: BRI, SME, 3DS, LDN.}
  \\ Voce aggregata su quattro video trascritti e letti per intero: la dissezione di un salvataggio di Rosso, il cavo Link negli emulatori, lo scambio locale su Switch in Rosso Fuoco, e due dimostrazioni dello strumento di trasferimento di riferimento. La trascrizione è avvenuta con lo strumento a riga di comando che scarica i sottotitoli automatici, ripuliti dal programma del progetto perché i sottotitoli a scorrimento ripetono ogni riga.

\bibitem{gbatemp-smeraldo}
  \textbf{GBAtemp, due discussioni sulla scrittura del salvataggio}.
  \\ \url{https://gbatemp.net/threads/save-failed-on-real-pokemon-emerald.645336/}
  \\ \emph{Livello 5, forum o community; letta. Track: SME.}
  \\ Voce aggregata su due discussioni lette dalle schermate consegnate dall'utente, perché il dominio respinge il recupero automatico. La prima riguarda il messaggio di salvataggio fallito su cartuccia genuina, la seconda la scrittura di un salvataggio da centoventotto kibibyte su una cartuccia di riproduzione la cui memoria è di sessantaquattro.

\bibitem{retroreversing}
  \textbf{RetroReversing, Game Boy}.
  \\ \url{https://www.retroreversing.com/gameboy}
  \\ \emph{Livello 4, articolo, blog o ricerca applicata; letta. Track: BRI.}
  \\ Raccolta di risorse di reverse engineering sulla piattaforma, letta per individuare strumenti e documentazione che il progetto non conosceva.

\bibitem{hackaday-ponte}
  \textbf{Hackaday, il ponte impossibile}.
  \\ \url{https://hackaday.com/2021/12/07/bridging-game-worlds-with-the-impossible-pokemon-trade/}
  \\ \emph{Livello 4, articolo, blog o ricerca applicata; letta. Track: BRI.}
  \\ Articolo divulgativo sul primo ponte fra generazioni, letto per estrarne la descrizione dell'hardware.

\bibitem{pokerom-trader}
  \textbf{pokerom-trader}.
  \\ \url{https://github.com/savaughn/pokerom-trader}
  \\ \emph{Livello 3, implementazione di riferimento; letta. Track: BRI.}
  \\ Strumento che esegue uno scambio fra due file di salvataggio delle prime due generazioni su calcolatore, scritto in C con una libreria di manipolazione dei salvataggi e un'interfaccia grafica propria, con ricalcolo dei checksum e regole di evoluzione da scambio.

\bibitem{video-distribuzioni}
  \textbf{Goppier, la ricreazione delle distribuzioni Gen 3}.
  \\ \url{https://www.youtube.com/watch?v=NKBb-YS34wg}
  \\ \emph{Livello 4, articolo, blog o ricerca applicata; letta. Track: EVT, BRI.}
  \\ Racconto tecnico della ricreazione di tutte le distribuzioni di eventi di terza generazione, inglesi, giapponesi e da GameCube, trascritto e letto per intero il 2026-08-28. Descrive come si trova il multiboot dentro la ROM di distribuzione confrontando i byte che passano sul cavo, come lo si decomprime attraverso la chiamata di sistema del BIOS, come si aggira il checksum additivo che ne difende l'integrità, e come si impostano i parametri dell'esemplare riusando il codice del gioco per indice di parametro.

\bibitem{video-ereader-eventi}
  \textbf{im a blisy, gli eventi e-Reader su cartuccia vera}.
  \\ \url{https://www.youtube.com/watch?v=fK-Actf6kME}
  \\ \emph{Livello 4, articolo, blog o ricerca applicata; letta. Track: EVT, SME.}
  \\ Guida esaustiva alle vie con cui un evento entra in una cartuccia originale, trascritta e letta per intero il 2026-08-28: l'emulatore con due istanze collegate e il BIOS reale, l'e-Reader con un salvataggio precostituito, la scheda riprogrammabile, la stampa delle carte con il codice a punti, e l'iniezione diretta della carta meraviglia nel salvataggio estratto.

\bibitem{video-pcny}
  \textbf{Hard4Games, la macchina del Pokemon Center di New York}.
  \\ \url{https://www.youtube.com/watch?v=AVhqlol6k9o}
  \\ \emph{Livello 4, articolo, blog o ricerca applicata; letta. Track: EVT.}
  \\ Preservazione dell'hardware e del software di distribuzione del Pokemon Center di New York, trascritta e letta per intero il 2026-08-28. Documenta i due dischi in formato proprietario di sviluppo, le schede di campagna che dichiarano quali esemplari distribuire e in quale finestra temporale, le schede che funzionano da chiave, e le date delle campagne di seconda e terza generazione.

\bibitem{video-gbi-bonus}
  \textbf{SuperrSonic, i contenuti bonus dal Game Boy Player}.
  \\ \url{https://www.youtube.com/watch?v=GBEMP2kEpPw}
  \\ \emph{Livello 4, articolo, blog o ricerca applicata; letta. Track: EVT.}
  \\ Programma per GameCube che applica a un salvataggio di cartuccia i contenuti bonus di diversi giochi, fra cui quelli dei giochi Pokemon di terza generazione, trascritto e letto per intero il 2026-08-28. La via passa dall'interfaccia alternativa del Game Boy Player, che invia un programma multiboot senza cavo perché il collegamento è interno alla console, ed è capace di estrarre BIOS, ROM e salvataggio e di ripristinare il salvataggio modificato.

\bibitem{bank-fine-servizio}
  \textbf{Nintendo, fine del servizio di Pokemon Bank}.
  \\ \url{https://en-americas-support.nintendo.com/app/answers/detail/a_id/61543}
  \\ \emph{Livello 2, wiki o riferimento di dominio; letta. Track: EVT, 3DS.}
  \\ Comunicazione ufficiale del supporto Nintendo sulla chiusura di Pokemon Bank, letta il 2026-08-28: il servizio termina giovedì 25 febbraio 2027 alle 19:00 PST, cioè il 26 febbraio alle 12:00 JST, e con esso cessa la possibilità di trasferire verso Pokemon Home. Chi ha esemplari depositati deve spostarli in Home prima di quella data, e la pagina non annuncia alcun periodo di tolleranza.

\bibitem{bulbapedia-trasferimenti}
  \textbf{Bulbapedia, la catena dei trasferimenti fra generazioni}.
  \\ \url{https://bulbapedia.bulbagarden.net/wiki/Pok\%C3\%A9mon_Bank}
  \\ \emph{Livello 2, wiki o riferimento di dominio; letta. Track: EVT, 3DS.}
  \\ Le tre pagine che descrivono la catena di trasferimento verso le piattaforme moderne, lette il 2026-08-28: Pokemon Bank con i giochi che Poke Transporter accetta come sorgente, cioè la quinta generazione e le riedizioni su Virtual Console della prima e della seconda; il Parco Amico, che porta dalla terza alla quarta generazione; e il Trasferitore, che porta dalla quarta alla quinta.

\bibitem{pkhex-eventi-gen3}
  \textbf{PKHeX, tabella degli eventi Gen 3 e vocabolario dei metodi}.
  \\ \url{https://github.com/kwsch/PKHeX}
  \\ \emph{Livello 3, implementazione di riferimento; letta. Track: EVT, BRI, SME.}
  \\ Le due parti di PKHeX che documentano le distribuzioni di terza generazione: la tabella `PKHeX.Core/Legality/Encounters/Data/Gen3/EncountersWC3.cs`, con centosettantasette voci ciascuna corredata del proprio metodo di generazione pseudocasuale, e l'enumerazione `PKHeX.Core/Legality/RNG/PIDType.cs`, che documenta i metodi uno per uno. La tabella vive nel codice e non in un documento perché, come dichiara il suo stesso commento, i dati di quella generazione non sono mai stati conservati in forma binaria uniforme e sono quindi scritti a mano.

\bibitem{bulbapedia-eventi-italia}
  \textbf{Bulbapedia, distribuzioni italiane di eventi in Gen 3}.
  \\ \url{https://bulbapedia.bulbagarden.net/wiki/List_of_Italian_event_Pok\%C3\%A9mon_distributions_in_Generation_III}
  \\ \emph{Livello 2, wiki o riferimento di dominio; letta. Track: EVT.}
  \\ Elenco delle distribuzioni italiane della terza generazione con luogo, date e campi di ciascun esemplare, letto il 2026-08-29. Porta la manifestazione del 2006 in un parco di divertimenti, dal ventitré al venticinque giugno, con allenatore di provenienza 10ANNI, identificativo 06227, dieci specie al livello settanta e la facoltà per il partecipante di scegliere tre esemplari, e la manifestazione del 2007 nel medesimo luogo, con un Mew di allenatore Aura e identificativo 20078.

\bibitem{gen-iii-event-patcher}
  \textbf{gen-iii-event-patcher}.
  \\ \url{https://github.com/superguideguy/gen-iii-event-patcher}
  \\ \emph{Livello 3, implementazione di riferimento; letta. Track: EVT, BRI.}
  \\ Strumento in Java che trasforma la ROM di un gioco di terza generazione in una ROM di distribuzione, e che applica uno script di evento a un salvataggio per uso personale. Letto il 2026-08-29 nella descrizione e nella struttura del sorgente, che comprende un compilatore elementare per gli script di evento e un costruttore di checksum.

\bibitem{project-wonder}
  \textbf{Project Wonder, distribuzioni per applicazione di differenze}.
  \\ \url{https://github.com/Goppier/Gen3DistributionRoms}
  \\ \emph{Livello 3, implementazione di riferimento; letta. Track: EVT.}
  \\ Raccolta di dieci differenze da applicare a una ROM di distribuzione preservata, letta il 2026-08-29. Produce distribuzioni degli oggetti di evento, cioè i biglietti, la mappa marina e la grotta alterna, più due eventi costruiti dalla comunità e tre distribuzioni di uova che riproducono quelle dei negozi, con supporto dichiarato a sei lingue fra cui l'italiano.

\bibitem{eventsgallery}
  \textbf{Project Pokemon, Events Gallery}.
  \\ \url{https://github.com/projectpokemon/EventsGallery}
  \\ \emph{Livello 3, implementazione di riferimento; catalogata e non letta. Track: EVT.}
  \\ Archivio collettivo della conservazione delle informazioni sugli eventi, per tutte le generazioni. Per la terza generazione le carte meraviglia si manipolano con lo strumento dedicato che questo registro già elenca.

\bibitem{pp-algoritmi-eventi}
  \textbf{Project Pokemon, ricerca sull'algoritmo di generazione degli eventi Gen 3}.
  \\ \url{https://projectpokemon.org/home/forums/topic/39517-gen-3-event-generation-algorithm-research-10anniv-etc/}
  \\ \emph{Livello 5, forum o community; letta. Track: EVT.}
  \\ Discussione sulla ricostruzione dell'algoritmo di generazione degli esemplari da evento della terza generazione, letta il 2026-08-29. Documenta che la determinazione dei metodi è avvenuta per reverse engineering a partire da campioni raccolti e non da una specifica, e riporta che il metodo di uno degli eventi dei negozi è stato risolto come variante a semi non ristretti in cui la metà alta del valore di personalità è messa in or esclusivo con la metà bassa, con l'identificativo dell'allenatore e con quello segreto.

\bibitem{frlg-home}
  \textbf{Pokemon, il collegamento della riedizione con il deposito in rete}.
  \\ \url{https://www.pokemon.com/us/news/pokemon-firered-version-and-pokemon-leafgreen-version-link-with-pokemon-home}
  \\ \emph{Livello 1, disassemblato o documentazione dell'hardware; letta. Track: ACE, EVT, LDN, 3DS.}
  \\ Annuncio del 13 agosto 2026, letto il 2026-08-31: le riedizioni per la console corrente dei due titoli della terza generazione si collegheranno al deposito in rete a ottobre 2026, con la versione 4.1.0 del servizio. Il trasferimento è a senso unico, la capienza del piano a pagamento sale da seimila a novemila esemplari, e il completamento dell'inventario consegna un esemplare altrimenti non ottenibile. È la fonte che corregge la scadenza di questo lavoro: per la sola terza generazione la chiusura del servizio di deposito cessa di essere l'ultima porta.

\bibitem{tpc-dati-alterati}
  \textbf{Pokemon, gestione dei dati alterati per mezzi non autorizzati}.
  \\ \url{https://support.pokemon.com/hc/en-us/articles/360055828671-Addressing-the-use-of-data-altered-via-unauthorized-means}
  \\ \emph{Livello 1, disassemblato o documentazione dell'hardware; letta. Track: ACE, GEN, EVT.}
  \\ Politica ufficiale, letta il 2026-08-31, ed è l'unico elemento non congetturale sul rischio delle vie che producono esemplari: chi risulti impiegare dati alterati può subire la restrizione del gioco in rete, la restrizione delle funzioni di scambio nella versione per dispositivo mobile e la sospensione dell'accesso al deposito, in forma temporanea o indefinita a discrezione del titolare e senza rimborso. Dichiara un'eccezione che delimita il rischio e non lo assolve: nessuna restrizione per chi possieda dati alterati senza intenzione, per esempio ricevendoli in uno scambio senza saperlo.

\bibitem{mankeymite-home}
  \textbf{MankeyMite, avvertenza sui trasferimenti verso il deposito in rete}.
  \\ \url{https://www.youtube.com/watch?v=KtJGkd0Qvvg}
  \\ \emph{Livello 4, articolo, blog o ricerca applicata; letta. Track: ACE, EVT, 3DS, LDN.}
  \\ Trascritta con yt-dlp e letta il 2026-08-31. Porta tre contributi: la catena ufficiale passo per passo, con i dettagli operativi sul luogo e sui prerequisiti di ciascun passaggio; il fatto che il trasferimento dal servizio di deposito verso il deposito in rete richieda il piano a pagamento di quest'ultimo, mentre il piano gratuito conserva trenta esemplari; e l'affermazione centrale sulla provenienza, cioè che il deposito in rete conserva sul proprio lato l'informazione della via da cui un esemplare è entrato, con la conseguenza che a parità di dati la storia differisce.

\bibitem{gen3-ace-builder}
  \textbf{Gen 3 ACE Pokemon Builder}.
  \\ \url{https://mankeymite.github.io/Gen3ACEPokemonBuilder/}
  \\ \emph{Livello 3, implementazione di riferimento; catalogata e non letta. Track: ACE, EVT.}
  \\ Funzione dichiarata dal suo autore: comporre il dato completo di un esemplare di terza generazione, comprese le vecchie distribuzioni di evento, con informazioni dell'allenatore, statistiche e lucentezza a scelta, restituendo un codice da digitare nei nomi delle scatole ed eseguire con il difetto del motore di testo. L'autore dichiara le opzioni di generazione conforme attive per difetto e dichiara di non poter garantire l'accettazione da parte del deposito in rete.

\bibitem{ace-archive}
  \textbf{Gen 3 ACE Archive}.
  \\ \url{https://mankeymite.github.io/gen3-ace-archive/}
  \\ \emph{Livello 3, implementazione di riferimento; letta. Track: ACE.}
  \\ Archivio di codici cercabili per titolo, lingua e piattaforma, con guide video, generatori e documentazione di allestimento, letto il 2026-08-31. Va registrata un'assenza perché è significativa quanto un contenuto: non contiene alcuna dichiarazione sulla legittimità degli esemplari prodotti, sulla loro accettazione da parte dei verificatori né sui rischi per l'account, e tratta la tecnica come problema di implementazione e non di conseguenze.

\bibitem{mankeymite-server}
  \textbf{MankeyMite's server, comunità dell'esecuzione di codice in generazione III}.
  \\ \url{https://discord.com/invite/rjQGPhG7e3}
  \\ \emph{Livello 5, forum o community; letta. Track: ACE, EVT, GEN, LDN.}
  \\ Ventiquattro canali e centounomilaottocentodiciassette messaggi, esportati il 2026-09-01 e interrogati per le domande dichiarate aperte anziché letti; le affermazioni citate sono attribuite al loro autore. È la fonte che corregge la formulazione della domanda sulla legittimità, stabilendo che i verificatori sono tre con severità decrescente. Porta inoltre il tracciatore univoco assegnato a ogni esemplare in ingresso nel deposito, la via composta attribuita al suo autore e in forma più semplice di quella ipotizzata da questo lavoro, il funzionamento dell'esecuzione di codice sulla riedizione per console moderna con i codici differenziati dal filtro sulle parole vietate, l'ammissione che l'elenco di quel filtro impiegato dal costruttore è una ricostruzione empirica, e la cifra di quattrocento esemplari da distribuzione con la dichiarazione che non si termina prima della chiusura del servizio.

\bibitem{insidegadgets-canale}
  \textbf{insideGadgets, canale di assistenza del produttore del lettore}.
  \\ \url{https://shop.insidegadgets.com/product/gbxcart-rw/}
  \\ \emph{Livello 5, forum o community; letta. Track: BAT, SME, BRI.}
  \\ Canale di assistenza del produttore del lettore di cartucce, esportato il 2026-08-31 con l'esportatore della comunità e letto per filtri il 2026-09-01: cinquantaquattromilasettecentocinquantuno messaggi, dai quali sono state estratte per parola chiave le testimonianze citate in questo lavoro, ciascuna attribuita al proprio autore. È la fonte di ciò che nessuna documentazione di prodotto contiene: la soglia di ritenzione dei componenti impiegati, la prova non strumentale per stabilire se la pila tenga ancora, il fatto che l'avvio dell'interfaccia nella modalità della terza generazione cancelli il salvataggio di una cartuccia della seconda anche senza comando di connessione, il formato della pila con la necessità della variante a linguette, e la sequenza in tre passi per rimettere in ordine l'orologio dopo la sostituzione.

\bibitem{home-checklist}
  \textbf{Pokemon HOME Checklist, inventario vivente e marchi di origine}.
  \\ \url{https://jacs720.github.io/Home-Checklist/}
  \\ \emph{Livello 3, implementazione di riferimento; letta. Track: ACE, GEN, EVT.}
  \\ La pagina è un'applicazione a pagina singola e una richiesta restituisce il solo involucro; il contenuto è stato ricavato il 2026-08-31 sondando per costanti il fascio JavaScript, da cui le enumerazioni e le etichette di interfaccia sono leggibili in chiaro anche dopo la minimizzazione. Ne viene il modello dei dati: undici marchi di origine, di cui uno dedicato alla terza generazione e trattato nel codice come non ancora ottenibile; undici collezioni di provenienza speciale; quindici profili di collezione predefiniti, la cui cardinalità differisce di un ordine di grandezza fra il minimo e quello per marchio; quattro livelli di reperibilità di cui uno dichiaratamente ipotetico. Il catalogo delle singole specie non compare in chiaro e non è stato enumerato, quindi la fonte documenta la tassonomia e non l'elenco.

\bibitem{berichandev}
  \textbf{berichandev, canale di trasmissione sui bot di scambio}.
  \\ \url{https://www.twitch.tv/berichandev}
  \\ \emph{Livello 5, forum o community; letta. Track: GEN.}
  \\ Registrato il 2026-08-31 su indicazione dell'utente, che lo descrive come noto nella comunità per avere sviluppato e ospitato bot di scambio automatico sui titoli per la console corrente; letto il 2026-09-01 su consegna di una schermata dei pannelli descrittivi del canale, che è la via della consegna manuale prevista dalla regola sulle fonti non recuperabili.

\bibitem{amiibodoctor-generazione}
  \textbf{amiibodoctor, generare esemplari per il titolo competitivo}.
  \\ \url{https://amiibodoctor.com/2025/08/25/how-to-generate-pokemon-to-use-in-pokemon-champions/}
  \\ \emph{Livello 4, articolo, blog o ricerca applicata; letta. Track: GEN.}
  \\ Registrata il 2026-08-31 e letta per intero il 2026-09-01, recuperata con una richiesta locale al primo tentativo dopo essere stata catalogata come non letta. È dell'agosto 2025 e in questo dominio un anno è molto, quindi va letta sapendo che la procedura descritta può essere stata resa obsoleta da un aggiornamento del titolo.

\bibitem{home-tracker-psa}
  \textbf{Il tracciatore del deposito e la regola che governa chi entra senza}.
  \\ \url{https://www.reddit.com/r/PokemonHome/comments/1vqrxf4/psa_pok\%C3\%A9mon_home_bank_connectivity_the_home/}
  \\ \emph{Livello 5, forum o community; letta. Track: PKD, EVT, ACE.}
  \\ Nota tecnica della comunità, letta il 2026-09-09 dal cluster del corpus sull'accesso alla banca, che enuncia in una riga la regola di ammissione del deposito: un esemplare privo di tracciatore viene accettato soltanto se il gioco che dichiara come origine coincide con il gioco da cui lo si sta depositando. Porta l'esempio e il controesempio, cioè un leggendario con marchio di origine della nona generazione accettato dentro un salvataggio di nona e il medesimo leggendario con origine di quarta rifiutato dal medesimo salvataggio. Documenta inoltre l'unica eccezione nota, cioè che gli esemplari depositati dal titolo dell'ottava generazione sono accettati qualunque sia la loro origine, e ne spiega il meccanismo con la forma dei dati, poiché quel titolo è il formato di sbarco che la banca impiega.

\bibitem{guide-origine-gen1}
  \textbf{Guida al catalogo per regione d'origine, prima generazione}.
  \\ \url{https://www.reddit.com/r/PokemonHome/comments/1p0hzxn/original_region_guides_gen_1/}
  \\ \emph{Livello 5, forum o community; letta. Track: PKD, EVT.}
  \\ Guida della comunità al completamento del catalogo di prima generazione con esemplari nativi di quella regione, letta il 2026-09-09. Dichiara che un esemplare selvatico di prima generazione non può essere legittimamente cromatico, mentre lo possono essere i doni, gli scambi in gioco e gli incontri fissi; enumera gli esclusivi di versione fra le edizioni giapponesi e internazionali, e le tredici specie che la terza versione non dà; elenca i nove scambi in gioco della coppia originale e i sette della terza versione con il loro soprannome.

\bibitem{guide-origine-gen2}
  \textbf{Guida al catalogo per regione d'origine, seconda generazione}.
  \\ \url{https://www.reddit.com/r/PokemonHome/comments/1p362dn/original_region_guides_gen_2/}
  \\ \emph{Livello 5, forum o community; letta. Track: PKD, EVT.}
  \\ Guida della comunità al completamento del catalogo di seconda generazione, letta il 2026-09-09. Enuncia le specie che quella generazione non dà legittimamente, cioè gli iniziali della prima, i fossili tranne uno, i tre uccelli leggendari e i due mitici; dichiara che Unown vi esiste nelle sole ventisei lettere, e che le due forme di punteggiatura compaiono per la prima volta nelle riedizioni di terza generazione; registra che il mitico dell'erba è ottenibile soltanto nella terza versione, che è anche la sola in cui lo si possa cercare cromatico.

\bibitem{guide-origine-gen3}
  \textbf{Guida al catalogo per regione d'origine, terza generazione}.
  \\ \url{https://www.reddit.com/r/PokemonHome/comments/1vphi1w/original_region_guides_gen_3/}
  \\ \emph{Livello 5, forum o community; letta. Track: PKD, EVT.}
  \\ Guida della comunità aggiornata al 15 agosto 2026, cioè dopo l'annuncio della chiusura, letta il 2026-09-09. Dichiara che la versione per console corrente dei due titoli di Kanto include l'evento dell'isola della nascita, che il mitico dello spazio profondo vi si può cercare cromatico e che registra Kanto come regione. Completa l'asse delle uova del deposito per console fissa con la mossa esclusiva di ciascuna delle quattro, e dichiara che soltanto una di esse deve dimenticarla per essere trasferita. Sui due giochi da console fissa chiarisce che non portano specie esclusive ma fiocchi e luoghi d'incontro esclusivi, e che ogni esemplare vi si può cercare cromatico a probabilità piena.

\bibitem{bulbapedia-differenze-sesso}
  \textbf{Bulbapedia, elenco delle specie con differenze di sesso visibili}.
  \\ \url{https://bulbapedia.bulbagarden.net/wiki/List_of_Pok\%C3\%A9mon_with_gender_differences}
  \\ \emph{Livello 2, wiki o riferimento di dominio; letta. Track: PKD.}
  \\ Pagina enciclopedica letta il 2026-09-08 con una richiesta locale e salvata in `\_notes/fonti/`. Enumera specie per specie le differenze visibili fra i due sessi, con il numero di catalogo a quattro cifre davanti a ciascuna riga, il che la rende estraibile a macchina; dichiara centodue specie nella prosa delle proprie sezioni e ne porta centodue nelle tabelle. Aggiunge fuori tabella la forma regionale di una specie che porta la differenza anche in quella forma.

\bibitem{bulbapedia-id-notevoli}
  \textbf{Bulbapedia, elenco degli identificativi di allenatore notevoli}.
  \\ \url{https://bulbapedia.bulbagarden.net/wiki/List_of_notable_ID_numbers}
  \\ \emph{Livello 2, wiki o riferimento di dominio; letta. Track: PKD, EVT.}
  \\ Tabella enciclopedica di ottocentonovantasei righe utili, ciascuna con l'identificativo, il nome dell'allenatore di provenienza e l'evento a cui appartiene, letta il 2026-09-09 con una richiesta locale dopo che l'attraversamento del corpus non l'aveva scaricata. Dichiara essa stessa tre righe come non appartenenti ad alcun esemplare da collezione, cioè gli identificativi degli esemplari da noleggio e quello che il catalogo di un gioco usa internamente.

\bibitem{sfide-impossibili}
  \textbf{Le sfide del deposito che la chiusura rende impossibili}.
  \\ \url{https://www.reddit.com/r/PokemonHome/comments/1voa411/all_pokemon_home_challenges_that_will_become/}
  \\ \emph{Livello 5, forum o community; letta. Track: PKD.}
  \\ Post della comunità che elenca trentasette sfide esplicite più cinque già invisibili fra quelle che la chiusura della banca renderà impossibili, letto il 2026-09-08. L'autore dichiara di averle ricavate dall'elenco di un archivio di terze parti, e nei commenti dichiara di avere voluto escludere le sfide che le riedizioni per console corrente rimetteranno in gioco, dimenticandone una.

\bibitem{bank-exclusives}
  \textbf{Raccolta di ciò che si perde con la chiusura della banca}.
  \\ \url{https://www.reddit.com/r/PokemonHome/comments/1apiuee/pok\%C3\%A9mon_bank_exclusives_masterpost/}
  \\ \emph{Livello 5, forum o community; letta. Track: PKD, EVT.}
  \\ Raccoglitore mantenuto da una collezionista, aggiornato al febbraio 2024 e letto il 2026-09-08 come voce principale del cluster sulla chiusura. Elenca sessantatre mosse fra tagliate, esclusive dei mitici e consegnate solo da distribuzioni, più tre abilità nella stessa condizione; elenca i fiocchi che ciascuna generazione ha conferito e le successive no; e nomina le combinazioni fra specie e sfera che nessun gioco moderno può produrre.

\bibitem{foglio-scambi-doni}
  \textbf{Cartella di calcolo degli scambi in gioco, dei doni e delle uova}.
  \\ \url{https://docs.google.com/spreadsheets/d/e/2PACX-1vTusrkjjQVxfqANVboZbw-VplOUBioFRDcHi5yEz4tXupmNvs9s2MHDBPA0jBzS38Ic9UT6Xulr0Sko/pubhtml}
  \\ \emph{Livello 5, forum o community; letta. Track: PKD, EVT.}
  \\ Cartella di calcolo comunitaria di ottocentouno voci fra scambi in gioco, doni, uova ed esemplari con cui si interagisce, scaricata il 2026-09-08 con una richiesta locale all'indirizzo di pubblicazione e divisa in quattro schede secondo l'epoca e la natura del titolo. Le due schede piccole, cioè quelle sui giochi derivati e sui giochi di deposito, sono l'elenco delle vie che non esistono più.

\bibitem{foglio-eventi-interni}
  \textbf{Cartella di calcolo degli esemplari da evento interno al gioco}.
  \\ \url{https://www.reddit.com/r/pokemontrades/comments/18y648l/}
  \\ \emph{Livello 5, forum o community; letta. Track: PKD.}
  \\ Cartella di calcolo comunitaria di millequattrocentocinquantotto esemplari consegnati o imposti una volta sola dalla storia dei giochi, più duecentotrentotto scambi in gioco su una scheda a parte, scaricata il 2026-09-08. La distribuzione per titolo, misurata sulla scheda esportata, dà novecentottantasei voci da titoli la cui via verso il deposito passa dalla banca e quattrocentosettantadue da titoli a via diretta.

\bibitem{sottolivellati}
  \textbf{Guida agli esemplari sottolivellati}.
  \\ \url{https://www.reddit.com/r/pokemon/comments/hhlm8k/}
  \\ \emph{Livello 5, forum o community; letta. Track: PKD.}
  \\ Guida comunitaria letta il 2026-09-08, da cui discende la scheda dei sottolivellati della lista di controllo che l'utente ha su disco e che la accredita per nome. Un esemplare sottolivellato è ottenuto a un livello inferiore a quello in cui la sua specie evolve per natura, quindi porta una proprietà che nessun altro titolo può ricreare; la guida dichiara le dieci voci più difficili con la loro probabilità, fra cui una su diecimila per un caso di quinta generazione e una su millecento per uno di terza.

\bibitem{timovm-ace-gen2}
  \textbf{Guide all'esecuzione di codice in seconda generazione}.
  \\ \url{https://glitchcity.wiki/wiki/Guides:TimoVM\%27s_gen_2_ACE_setups}
  \\ \emph{Livello 2, wiki o riferimento di dominio; letta. Track: ACE, PKD, EVT.}
  \\ Guide della wiki dedicata ai difetti sfruttabili, lette il 2026-09-09, che coprono ogni versione e ogni lingua dei tre titoli di seconda generazione, l'italiano compreso, e sono dichiarate compatibili sia con le cartucce originali sia con le riedizioni per console corrente. L'allestimento si raggiunge dopo la seconda medaglia nei due titoli maggiori e senza alcuna medaglia nel terzo, e la fonte dichiara fra trenta e sessanta minuti per installare il programma di cinquanta byte che consente di scrivere ed eseguire codice arbitrario senza più passare dalle quantità degli oggetti.

\bibitem{glitchcity-mailwriter}
  \textbf{Glitch City, i codici del programma di scrittura}.
  \\ \url{https://glitchcity.wiki/wiki/Guides:Mail_Writer_Codes}
  \\ \emph{Livello 2, wiki o riferimento di dominio; letta. Track: ACE, PKD.}
  \\ Raccolta dei codici già scritti per il programma di cinquanta byte installato dalle guide di seconda generazione, letta il 2026-09-09. Tre di essi riguardano la produzione: uno consegna un mitico al livello cinque e ne adatta nome e identificativo dell'allenatore perché il trasferimento lo accetti, uno fa lo stesso per la linea di due iniziali di terza generazione, e uno riempie la scatola attiva chiedendo ogni volta specie e livello dell'esemplare successivo. Accanto stanno codici che cambiano le condizioni di una caccia invece di produrre, e un programma più grande con cui si leggono e si cambiano i valori in memoria.

\bibitem{pkmnclassic}
  \textbf{Rete ricostruita dei servizi in rete di quarta e quinta generazione}.
  \\ \url{https://pkmnclassic.net/}
  \\ \emph{Livello 3, implementazione di riferimento; letta. Track: EVT, PKD.}
  \\ Servizio non ufficiale che fornisce oggi il sistema di scambio globale, i video delle lotte e altri servizi in rete per la quarta e la quinta generazione, raggiungibile cambiando un solo parametro sulla console. La pagina è stata letta il 2026-09-08 nel momento in cui la si scarica e dichiara contatori di esercizio aggiornati, cioè esemplari offerti e video registrati, il che è la prova che il servizio è vivo e non una pagina abbandonata.

\bibitem{shacknews-dns}
  \textbf{Come riaprire i doni segreti di quarta e quinta generazione}.
  \\ \url{https://www.shacknews.com/article/108512/how-to-unlock-gen-4-and-5-pokemon-mystery-events-in-2018}
  \\ \emph{Livello 4, articolo, blog o ricerca applicata; letta. Track: EVT, PKD.}
  \\ Articolo redazionale del 2018 che descrive la procedura per riaprire i doni segreti nei giochi di quarta e quinta generazione cambiando il nome del server dei nomi sulla console, e che dichiara il sistema di scambio globale pienamente funzionante in quarta generazione. La consegna dei doni vi è descritta come casuale e ripetibile, e produce duplicati.

\bibitem{sinjoh-ruins}
  \textbf{La procedura delle rovine e il terzetto al livello uno}.
  \\ \url{https://www.reddit.com/r/PokemonHGSS/comments/1hy3b12/}
  \\ \emph{Livello 5, forum o community; letta. Track: PKD, EVT.}
  \\ Guida passo per passo letta il 2026-09-09, che chiude un punto che il progetto teneva aperto: il terzetto del tempo e dello spazio al livello uno si ottiene portando un esemplare mitico di quarta generazione in una riedizione di seconda, tenendolo come solo esemplare in squadra, entrando nelle rovine e scegliendo uno dei tre cerchi, ciascuno dei quali dà una delle tre specie al livello uno.

\bibitem{bulbapedia-macchine-nascoste}
  \textbf{Bulbapedia, le macchine nascoste per generazione}.
  \\ \url{https://bulbapedia.bulbagarden.net/wiki/HM}
  \\ \emph{Livello 2, wiki o riferimento di dominio; letta. Track: PKD, EVT.}
  \\ Pagina enciclopedica che elenca quali mosse siano macchine nascoste in ciascuna generazione, consegnata il 2026-09-09 dentro una ricerca dell'utente e assorbita nel documento della catena. È la tavola su cui poggia il controllo automatico che misura quante voci dei lotti prodotti il passaggio fra generazioni rifiuterebbe.

\bibitem{surfing-pikachu}
  \textbf{La storia dell'esemplare surfista, dal titolare della serie}.
  \\ \url{https://www.pokemon.com/us/pokemon-news/waxing-nostalgic-about-surfing-pikachu}
  \\ \emph{Livello 2, wiki o riferimento di dominio; letta. Track: EVT, PKD.}
  \\ Articolo del sito ufficiale della serie, consegnato il 2026-09-09 dentro la medesima ricerca, che racconta la storia dell'esemplare consegnato come premio di un torneo per console fissa e della mossa che lo caratterizza. Accanto a esso una voce enciclopedica indipendente stabilisce che quella specie non apprende quella mossa per macchina né per insegnamento in alcun titolo successivo.

\bibitem{pokepc-dati}
  \textbf{PokePC Classic, i dati del tracciatore di catalogo vivente}.
  \\ \url{https://github.com/pokepc/classic.pokepc.net}
  \\ \emph{Livello 3, implementazione di riferimento; letta. Track: PKD.}
  \\ Il tracciatore gia' noto come SuperEffective.gg pubblica sotto licenza permissiva i dati che lo alimentano, letti il 2026-09-10 dal deposito pubblico e non dall'interfaccia: un'anagrafica di millecinquecentonovantanove voci con il numero di catalogo e i contrassegni che dicono che cosa ciascuna sia, e sette disposizioni in scatole per il deposito, che sono l'elenco delle caselle da riempire perche' un catalogo sia completo. Le sette non differiscono per il solo ordinamento come sembrerebbe: due danno una casella propria alla forma gigamax e due si dichiarano minime, quindi i totali sono millequattrocentoventicinque, milletrecentottantasette e milletrecentosettantatre.

\bibitem{monarium}
  \textbf{Monarium, generatore di disposizioni del deposito con marchi}.
  \\ \url{https://github.com/kristopheles/monarium}
  \\ \emph{Livello 3, implementazione di riferimento; letta. Track: PKD.}
  \\ Tracciatore pubblicato nel 2026 e letto il 2026-09-10 nella presentazione del suo autore, che va installato per conto proprio perche' non esiste un'istanza pubblica. Genera la disposizione delle scatole a partire da piu' di quindici opzioni, fra cui le sole evoluzioni finali, le varianti regionali accanto alla forma di partenza, i cromatici in scatole dedicate e l'esclusione delle forme da evento; sopra la disposizione tiene cinque contrassegni per esemplare, cioe' se porti l'allenatore di chi colleziona, se sia stato catturato nel titolo o nella regione di origine, se sia cromatico, se stia nella sfera giusta e se venga dall'applicazione per telefono.

\bibitem{foglio-arca}
  \textbf{Cartella di calcolo del catalogo per coppia di sessi}.
  \\ \url{https://www.reddit.com/r/PokemonHome/comments/1djbdm0/}
  \\ \emph{Livello 5, forum o community; letta. Track: PKD.}
  \\ Foglio comunitario del 2024, letto il 2026-09-10, che definisce operativamente il profilo detto dell'arca: un esemplare per ciascuno dei due sessi di ogni specie e di ogni forma, uno solo per le specie senza sesso o con un sesso solo, e la variante cromatica di ciascuno per chi voglia il caso estremo. L'autore dichiara di escludere gli esemplari a distribuzione limitatissima e di comprendere le sessanta e piu' configurazioni della specie il cui dolcetto non e' un campo della forma.

\bibitem{foglio-collezione-thundrosaur}
  \textbf{Cartella di calcolo della collezione, sei tracciatori in uno}.
  \\ \url{https://www.reddit.com/r/PokemonHome/comments/xp1ise/}
  \\ \emph{Livello 5, forum o community; letta. Track: PKD.}
  \\ Foglio comunitario del 2022, letto il 2026-09-10, che l'autore dichiara composto prendendo gli elementi migliori da piu' strumenti precedenti. Tiene insieme sei tracciatori distinti: il catalogo vivente, il catalogo degli esemplari di taglia massima, i marchi di origine, i marchi ordinari, i fiocchi e le sfere di cattura.

\bibitem{dns-eventi-gen45}
  \textbf{I doni di quarta e quinta generazione dal servizio ricostruito, elenco e testimonianze}.
  \\ \url{https://www.reddit.com/r/wiimmfi/comments/fugd8l/}
  \\ \emph{Livello 5, forum o community; letta. Track: PKD, EVT.}
  \\ Post del 2020 letto il 2026-09-10, con l'elenco dei doni che il suo autore dichiara di avere ricevuto cambiando i server dei nomi sulla console e collegandosi al servizio ricostruito, titolo per titolo, e con milletrecentottantaquattro commenti che lo aggiornano. L'elenco comprende per la quarta generazione i mitici, i tre cani cromatici, alcuni oggetti chiave e due doni che non furono mai distribuiti ufficialmente, cioe' il flauto azzurro e la capsula; per la quinta comprende ventidue voci fra cui i tre draghi cromatici e i mitici della generazione. Due testimonianze indipendenti del 2026-09-07, cioe' tre giorni prima della lettura, dichiarano il canale ancora attivo e nominano le sole voci che non hanno ottenuto.

\medskip
\noindent \textbf{Riferimenti teorici.} Le voci che seguono sono la letteratura canonica dei concetti impiegati nell'analisi quantitativa, citate per attribuzione del concetto e non come fonti consultate nel corso del lavoro: la definizione di entropia si attribuisce a Shannon perché è sua, non perché quell'articolo sia stato aperto qui. I numeri di pagina non sono riportati dove non sono stati verificati.\\[0.8ex]

\bibitem{shannon48}
  \textbf{C. E. Shannon}. A Mathematical Theory of Communication.
  \\ Bell System Technical Journal, vol. 27, pp. 379-423 e 623-656, 1948.
  \\ \emph{Riferimento teorico, citato per attribuzione.}
  \\ Definizione di entropia di una sorgente discreta, di informazione mutua e di capacità di canale. È il riferimento per ogni misura in bit impiegata in questo lavoro.

\bibitem{shannon49}
  \textbf{C. E. Shannon}. Communication Theory of Secrecy Systems.
  \\ Bell System Technical Journal, vol. 28, n. 4, pp. 656-715, 1949.
  \\ \emph{Riferimento teorico, citato per attribuzione.}
  \\ Nozione di sicurezza perfetta e dimostrazione che essa richiede una chiave di entropia almeno pari a quella del messaggio. È il criterio contro cui si misura la cifratura della generazione 3.

\bibitem{vernam26}
  \textbf{G. S. Vernam}. Cipher Printing Telegraph Systems for Secret Wire and Radio Telegraphic Communications.
  \\ Journal of the American Institute of Electrical Engineers, vol. 45, pp. 109-115, 1926.
  \\ \emph{Riferimento teorico, citato per attribuzione.}
  \\ Il cifrario a somma modulo due fra messaggio e chiave, di cui la cifratura della generazione 3 è un'istanza con chiave riusata.

\bibitem{cover-thomas}
  \textbf{T. M. Cover, J. A. Thomas}. Elements of Information Theory, 2ª edizione.
  \\ Wiley, 2006.
  \\ \emph{Riferimento teorico, citato per attribuzione.}
  \\ Trattazione sistematica di entropia, entropia condizionata e informazione mutua, e della disuguaglianza di elaborazione dei dati, impiegata per stabilire che la permutazione non aggiunge incertezza dato il valore di personalità.

\bibitem{katz-lindell}
  \textbf{J. Katz, Y. Lindell}. Introduction to Modern Cryptography, 3ª edizione.
  \\ CRC Press, 2020.
  \\ \emph{Riferimento teorico, citato per attribuzione.}
  \\ Enunciato e dimostrazione moderna del teorema sulla sicurezza perfetta, e trattazione dell'attacco per sovrapposizione su una chiave riusata.

\bibitem{peterson-brown}
  \textbf{W. W. Peterson, D. T. Brown}. Cyclic Codes for Error Detection.
  \\ Proceedings of the IRE, vol. 49, n. 1, pp. 228-235, 1961.
  \\ \emph{Riferimento teorico, citato per attribuzione.}
  \\ Introduzione dei codici ciclici per la rilevazione d'errore e delle loro garanzie sugli errori a raffica, che è il termine di confronto per il checksum additivo dei giochi.

\bibitem{lin-costello}
  \textbf{S. Lin, D. J. Costello}. Error Control Coding, 2ª edizione.
  \\ Pearson Prentice Hall, 2004.
  \\ \emph{Riferimento teorico, citato per attribuzione.}
  \\ Capacità di rilevazione di un codice in funzione della propria distanza minima, e limite della probabilità di errore non rilevato per un codice a sindrome corta.

\bibitem{rfc1071}
  \textbf{R. Braden, D. Borman, C. Partridge}. Computing the Internet Checksum, RFC 1071.
  \\ Internet Engineering Task Force, 1988.
  \\ \emph{Riferimento teorico, citato per attribuzione.}
  \\ La somma a complemento a uno con ripiegamento del riporto, cioè la tecnica che la generazione 3 impiega per ridurre a sedici bit una somma calcolata a trentadue.

\bibitem{rfc1662}
  \textbf{W. Simpson (a cura di)}. PPP in HDLC-like Framing, RFC 1662.
  \\ Internet Engineering Task Force, 1994.
  \\ \emph{Riferimento teorico, citato per attribuzione.}
  \\ La trasparenza a byte, cioè il byte stuffing con sequenza di fuga, di cui la lista di correzione del cavo Link è l'alternativa a lunghezza fissa.

\bibitem{rfc3927}
  \textbf{S. Cheshire, B. Aboba, E. Guttman}. Dynamic Configuration of IPv4 Link-Local Addresses, RFC 3927.
  \\ Internet Engineering Task Force, 2005.
  \\ \emph{Riferimento teorico, citato per attribuzione.}
  \\ Definizione del blocco di indirizzi link-local 169.254.0.0/16 e della sua autoconfigurazione, impiegato dal protocollo di rete locale della console.

\bibitem{tanenbaum}
  \textbf{A. S. Tanenbaum, D. J. Wetherall}. Computer Networks, 5ª edizione.
  \\ Prentice Hall, 2011.
  \\ \emph{Riferimento teorico, citato per attribuzione.}
  \\ Delimitazione di trama e trasparenza dei dati, rilevazione d'errore a livello di collegamento, e assegnazione dei canali nelle reti locali senza filo.

\bibitem{proakis}
  \textbf{J. G. Proakis, M. Salehi}. Digital Communications, 5ª edizione.
  \\ McGraw-Hill, 2008.
  \\ \emph{Riferimento teorico, citato per attribuzione.}
  \\ Modello del canale, occupazione di banda e relazione fra velocità di simbolo e larghezza spettrale, impiegati nel calcolo della non sovrapposizione dei canali.

\bibitem{gray-neuhoff}
  \textbf{R. M. Gray, D. L. Neuhoff}. Quantization.
  \\ IEEE Transactions on Information Theory, vol. 44, n. 6, pp. 2325-2383, 1998.
  \\ \emph{Riferimento teorico, citato per attribuzione.}
  \\ Teoria della quantizzazione scalare, quantizzatori non uniformi e saturazione, che è la forma esatta della conversione dell'allenamento fra le generazioni.

\bibitem{feller}
  \textbf{W. Feller}. An Introduction to Probability Theory and Its Applications, vol. 1, 3ª edizione.
  \\ Wiley, 1968.
  \\ \emph{Riferimento teorico, citato per attribuzione.}
  \\ Distribuzione geometrica con il proprio valore atteso e la propria varianza, e approssimazione di Poisson della binomiale, impiegate per il campionamento con rifiuto e per il conteggio delle occorrenze del byte riservato.

\bibitem{vonneumann51}
  \textbf{J. von Neumann}. Various Techniques Used in Connection with Random Digits.
  \\ National Bureau of Standards, Applied Mathematics Series, vol. 12, 1951.
  \\ \emph{Riferimento teorico, citato per attribuzione.}
  \\ Formulazione originale del metodo di campionamento con rifiuto, che è l'algoritmo con cui l'implementazione di riferimento genera il valore di personalità.

\bibitem{devroye}
  \textbf{L. Devroye}. Non-Uniform Random Variate Generation.
  \\ Springer-Verlag, 1986.
  \\ \emph{Riferimento teorico, citato per attribuzione.}
  \\ Trattazione sistematica del campionamento con rifiuto, compreso il costo atteso in funzione della probabilità di accettazione.

\bibitem{knuth2}
  \textbf{D. E. Knuth}. The Art of Computer Programming, vol. 2: Seminumerical Algorithms, 3ª edizione.
  \\ Addison-Wesley, 1997.
  \\ \emph{Riferimento teorico, citato per attribuzione.}
  \\ Aritmetica modulare e teorema cinese del resto, impiegato per stabilire l'indipendenza fra vincoli che agiscono su moduli coprimi.

\bibitem{knuth3}
  \textbf{D. E. Knuth}. The Art of Computer Programming, vol. 3: Sorting and Searching, 2ª edizione.
  \\ Addison-Wesley, 1998.
  \\ \emph{Riferimento teorico, citato per attribuzione.}
  \\ Rappresentazione delle permutazioni mediante tavole di inversione e sistema numerico fattoriale, che è la struttura della tabella di permutazione delle sottostrutture.

\bibitem{ieee80211}
  \textbf{IEEE}. IEEE Standard 802.11: Wireless LAN Medium Access Control and Physical Layer Specifications.
  \\ Institute of Electrical and Electronics Engineers, 2020.
  \\ \emph{Riferimento teorico, citato per attribuzione.}
  \\ Struttura dei frame di gestione e di azione, piano dei canali nella banda a 2,4 GHz e occupazione di banda delle portanti, su cui il protocollo di rete locale della console è costruito.

\end{thebibliography}
